\section*{\S~2e Mitgliederversammlung}
\label{sec:mitgliederversammlung}
\addcontentsline{toc}{section}{\S~2e Mitgliederversammlung}

\begin{enumerate}[label=(\arabic*)]
    \item Die Mitgliederversammlung besteht aus allen aktiven Mitgliedern, allen Mitgliedern der Reserve und allen Mitgliedern außer Dienst der 1st Aircat Division.
    \item Aktives Wahlrecht haben alle Mitglieder der 1st Aircat Division. Ein passives Wahlrecht haben nur die aktiven Mitglieder der 1st Aircat Division.
    \item \label{itm:zusammentretten_mitgliederversammlung} Die Mitgliederversammlung tritt bei Bedarf zusammen, mindestens aber einmal im Kalenderjahr. Damit die Mitgliederversammlung zusammentritt, muss die Leitung diese einberufen. Auf Verlangen von mindestens einem Drittel der Mitglieder muss die Leitung die Mitgliederversammlung einberufen. Sollte die Mitgliederversammlung auf einem öffentlichen Kanal der 1st Aircat Division zusammentreten, so muss diese mindestens 7 Tage im Voraus angekündigt werden. Sollten lediglich Abstimmungen und Wahlen über einen anderen Kanal durchgeführt werden, müssen diese für mindestens 7 Tage offen stehen. Anträge für die Mitgliederversammlung müssen spätestens 7 Tage vor der öffentlichen Einladung bei der Unit-Leitung eingegangen sein. Ansonsten ist die Einbringung eines Antrages nur noch während der Mitgliederversammlung möglich. Die Person, die einen Antrag einbringt, oder eine von ihr beauftragte dritte Person, muss anwesend sein. Für Anträge muss es während der Mitgliederversammlung die Möglichkeit zur Diskussion geben.
    \item Die Mitgliederversammlung tritt entweder an einem durch die Leitung angesetzten Termin auf einem öffentlichen Kanal der 1st Aircat Division zusammen, oder durch mehrere Abstimmungen über andere Kanäle, in denen über die zur Abstimmung stehenden Themen abgestimmt wird, beziehungsweise die nötigen Wahlen durchgeführt werden.
    \begin{enumerate}[label=(\Alph*)]
        \item \label{itm:regularer_Termin}Der Termin für die reguläre, jährliche Mitgliederversammlung muss einen Monat vor dem geplanten Termin von der Leitung dem erweiterten Stab vorgelegt und von diesem genehmigt werden. Existiert kein erweiterter Stab, so ist dieser Termin dem Stab vorzulegen. Der Termin muss dann zeitnah den Mitgliedern in geeigneter Art und Weise kommuniziert werden, sodass alle Fristen nach \S~2e Absatz \ref{itm:zusammentretten_mitgliederversammlung} eingehalten werden können. Der Termin sollte in den Dezember des jeweiligen Kalenderjahres und möglichst mindestens eine Woche vor Heiligabend liegen. Diese Mitgliederversammlung sollte in Präsenz durchgeführt werden.
    \end{enumerate}
    \item Ob eine Entscheidung durch die Mitgliederversammlung abgesegnet werden muss, entscheidet die Leitung. Alternativ können binnen 72 Stunden nach Bekanntgabe einer Entscheidung die Mitglieder der 1st Aircat Division eine Mitgliederversammlung verlangen. Eine Mitgliederversammlung findet dann statt, wenn mindestens ein Drittel der Mitglieder diesen Antrag unterstützen. Änderungen an dieser Satzung (ausgenommen redaktioneller Änderungen), die Wahl der Leitung oder die Beförderung der Leitung müssen durch die Mitgliederversammlung vorgenommen werden.
    \item Die Mitgliederversammlung ist beschlussfähig, wenn mindestens ein Drittel + 1 aller Mitglieder anwesend ist. Hierbei wird stets aufgerundet. Als anwesend gilt ein Mitglied, wenn es entweder physisch anwesend ist oder sich an einer Abstimmung beteiligt hat. Die Beschlussfähigkeit muss überprüft werden. Handelt es sich bei einem Drittel + 1 aller Mitglieder um passive Mitglieder, so ist die Mitgliederversammlung bei Anwesenheit einer Hälfte + 1 der aktiven Mitglieder beschlussfähig.
    \item\label{itm:satzungsaenderungen} Für Änderungen an dieser Satzung ist eine Zweidrittelmehrheit der abgegebenen Stimmen nötig. Redaktionelle Änderungen sind davon nicht betroffen. Bei allen anderen Entscheidungen und Beförderungen ist eine relative Mehrheit ausreichend. Satzungsänderungen treten sofort nach Bekanntgabe des Abstimmungsergebnisses in Kraft. Die neue Version der Satzung ist unverzüglich auf allen öffentlichen Kanälen der 1st Aircat Division zu veröffentlichen und alle Mitglieder sind darauf hinzuweisen.
    \item Anträge zur Tagesordnung einer Mitgliederversammlung sind unzulässig.
\end{enumerate}